\section{The Non-Closure of Individual Intelligence}


\subsection{Overview}


\S{}7 では、正しさの累積が分岐爆発を生み、個体内の知性的動態が停止条件を欠いた不安定化レジームに陥ることを動力学的に示した。本セクションはその構造的根拠を提示する。すなわち:

\begin{quote}
個体内には、知性的動態を停止させる演算子が \textbf{構造的に存在しえない}。
\end{quote}

この主張は心理学的観察ではない。意志の弱さ、知性の不足、訓練の欠如 ── これらいずれの問題でもない。それは、因果推論そのものの内部構造から導出される \textbf{構造的事実} である。

本セクションが論証するのは「個体は停止しない」ではなく「個体は \textbf{構造的に停止できない}」である。両者の違いは決定的である。前者は個体の選択や能力の問題として読みうるが、後者は構造そのものが要請する不可避の事実として読まれる。この区別が、\S{}9 で展開される越境境界の必然性の哲学的根拠となる。

なお、本セクションは \S{}5 で示した cross-domain 性のもとで論じられる。個体内非閉包は生物個体だけでなく、\S{}5.2 で扱った Transformer のような人工的個体系にも適用される。すなわち単一の Transformer モデルが訓練範囲を超えて知性的動態を停止させることができないという経験的事実は、本セクションが示す構造的事実の一つの現れである。

\subsection{The recursive structure of causal reasoning}


\subsubsection{The form of why-regress}


因果推論の特徴的構造は、次の再帰形式に集約される:

\begin{itemize}
\item $A$ はなぜ起こっているか?
\item $A$ は $B$ によって説明される
\item では $B$ はなぜ起こっているか?
\item $B$ は $C$ によって説明される
\item では $C$ はなぜ起こっているか...
\end{itemize}

この構造を \textbf{why-regress}(なぜの回帰)と呼ぶ。why-regress は推論者の悪癖ではない ── それは因果推論が因果推論であるための \textbf{構造的本質} である。

因果関係を問うこと自体が、答えが得られた瞬間に \textbf{新たな因果関係への問い} を構造的に生成する。これは選択ではなく構造である。$A$ が $B$ によって説明されるという答えを得たとき、その答えは:

\begin{itemize}
\item $B$ を新たな状態として \textbf{状態空間に導入} する(\S{}3.1 の連結状態空間に新ノードを追加)
\item $B$ を $A$ と関係づける \textbf{情報流 $J$ を生成} する(\S{}3.2 の時間順序付けに新エッジ)
\item $B$ の存在自体が \textbf{新たな問い} の起点となる(self-generating inference, \S{}7.2.2)
\end{itemize}

これらはすべて、答えを得ることの \textbf{構造的副産物} である。問いを止めようとして答えを得ても、答えは別の問いを生成する。

\subsubsection{The regress contains no intrinsic stopping point}


why-regress が構造的に停止しないことを、より厳密に論証する。

ある説明連鎖が時点 $t_n$ で「停止した」と判定されるためには、$t_n$ における説明状態が以下のいずれかを満たす必要がある:

\begin{itemize}
\item (a) その状態に対して、もはや「なぜ」を問う構造的根拠が存在しない
\item (b) その状態は自己説明的である(self-explanatory)
\item (c) その状態は最終的である(ultimate)と公理的に承認されている
\end{itemize}

しかし、因果推論の内部構造を観察すると、これらいずれも個体的推論からは \textbf{内発的に生成されない}:

(a) について:いかなる説明状態 $B$ も、その存在条件・適用範囲・例外を問う構造的根拠を内在的に持つ(\S{}7.2.1)。「もはや問えない」状態を個体的推論が自分で確定することはできない。

(b) について:「自己説明的」という属性は、その属性を判定する基準を必要とする。だが基準もまた説明を要する。自己説明性の判定は、自己説明性の基準への回帰を生むだけである。

(c) について:「最終的」という属性は、外的な公理化または承認を必要とする。個体的推論の内部からは、何かを「最終」と承認する根拠が構造的に得られない。

すなわち、why-regress の停止は \textbf{個体的推論の外側} からのみ生じうる。これが本セクションの中核的構造論的事実である。

\subsection{Conditions required for stable stopping}


推論が安定的に停止するための条件を整理する。これらは個体的推論からは内発的に生成されないが、外的に与えられれば停止を実現する条件である。

\subsubsection{External response}


第一の条件は \textbf{外部応答} である。他者または外部系から返される反応が、推論連鎖を停止させる構造的役割を果たす。

具体的な機構:

\begin{itemize}
\item 他者の応答は、推論連鎖の内部からは予測不可能な \textbf{新規情報} をもたらす
\item この新規情報は、推論連鎖の現在地点を \textbf{外部視点から評価} する材料となる
\item 外部視点からの評価は、内部からは生成できない「ここで停止する」という判断を提供する
\end{itemize}

外部応答が停止を可能にするのは、推論者が応答を「聞く」からではなく、応答が \textbf{推論連鎖の構造の外側から構造を制約する} からである。これは情報内容の問題ではなく、構造的位置の問題である。

\subsubsection{Shared agreement}


第二の条件は \textbf{共有された合意} である。複数の主体間で確立された前提・規範・参照点が、推論の境界を設定する。

合意の構造的役割:

\begin{itemize}
\item 合意は個別主体に \textbf{外在的} である(個体内では生成できない)
\item 合意は推論の \textbf{共通基盤} を提供する(議論の出発点を固定する)
\item 合意は逸脱に対する \textbf{構造的抵抗} を生む(個体的推論が共有基盤から離れすぎることを防ぐ)
\end{itemize}

ここで重要なのは、合意の \textbf{正しさ} ではなく \textbf{共有性} である。合意は誤っていても停止機能を果たしうる。停止を可能にするのは合意の真理値ではなく、その共有性が個体的推論の外側に立っているという構造的事実である。

\subsubsection{Externalization as log}


第三の条件は \textbf{log としての外化} である。推論過程が文字・記録・データとして個体外に固定されること。

log の構造的役割:

\begin{itemize}
\item log は推論過程を \textbf{時間外} に固定する(個体的時間流から切り離す)
\item log は推論を \textbf{複数主体に参照可能} にする(個体的閉包から解放する)
\item log は推論の \textbf{再開可能性} を保証する(停止と再開を分離する)
\end{itemize}

log が停止を可能にするのは、log が記憶の置換ではなく、推論連鎖を \textbf{個体的時間流の外側} に固定する機構だからである。log なしには、推論は停止しても保持されない ── 停止と消滅が区別されない。

\subsubsection{Temporal delay}


第四の条件は \textbf{時間的遅延} である。推論連鎖に意図的な間隔を導入すること。

時間的遅延の構造的役割:

\begin{itemize}
\item 遅延は推論連鎖の \textbf{自己生成的勢い} を断ち切る
\item 遅延は新たな情報が到来する \textbf{時間窓} を開く
\item 遅延は推論者自身が現在の連鎖を \textbf{外側から見る} 機会を提供する
\end{itemize}

時間的遅延は、上記三条件と比べると最も「内在的」に見えるが、構造的には個体的推論から内発的に生成されない。個体的推論は加速こそすれ、自発的に遅延を導入する構造的根拠を持たない。遅延を導入するには、推論連鎖の外側から「ここで止める」という介入が必要であり、これは個体内では生じない。

\subsubsection{The structural commonality}


これら四条件の共通する構造的特徴は、いずれも \textbf{個体的推論の外側} に位置していることである:

\begin{itemize}
\item 外部応答 = 他者という外側
\item 共有合意 = 共同体という外側
\item log = 時間と参照可能性という外側
\item 時間的遅延 = 推論連鎖の動的勢いの外側
\end{itemize}

四条件はすべて、推論を停止させるために \textbf{個体内では生成されない} 何かが必要であることを示している。

\subsection{The absence of an internal stopping operator}


\subsubsection{What the individual system can do}


個体系は以下を行うことができる:

\begin{itemize}
\item 推論を実行する(infer)
\item 説明を精緻化する(refine explanations)
\item 正しい知識を累積する(accumulate correct knowledge)
\item 推論連鎖を展開する(extend chains of reasoning)
\item 矛盾を検出し修正する(detect and resolve contradictions)
\end{itemize}

これらはいずれも、\S{}3 の三公理を満たす知性的動態の正当な作動である。

\subsubsection{What the individual system cannot do}


個体系が \textbf{安定的かつ自己持続的に} 行えないことは以下である:

\begin{itemize}
\item 「ここで停止する」と判断する(decide "here we stop")
\item 探究を崩壊させずに保留する(suspend inquiry without collapse)
\item 持続的な保持状態を維持する(maintain a persistent holding state)
\item 推論連鎖の終端を内発的に確定する(intrinsically determine the terminal of a chain)
\end{itemize}

これらはいずれも、停止演算子の \textbf{作動結果} として現れる動的状態である。停止演算子が存在しなければ、これらの状態は構造的に成立しない。

\subsubsection{The structural absence}


個体系における停止演算子の不在を、形式的にまとめる:

\begin{quote}
因果推論の構造を $\mathcal{R}$ とし、$\mathcal{R}$ に作用して停止を生成する演算子を $\mathcal{T}$ とする。
$\mathcal{T}$ は $\mathcal{R}$ の内部からは構造的に生成されない:
$$
\mathcal{T} \notin \text{Internal}(\mathcal{R})
$$
\end{quote}

この不在は、個体系の \textbf{欠陥} ではない。それは因果推論の構造そのものから導出される構造的事実である。$\mathcal{R}$ がより精緻になっても、より高速になっても、より多くの知識を含んでも、$\mathcal{T}$ は $\mathcal{R}$ の内部に現れない。

これは \S{}7.5 で示した「より多くの知性は解決にならない」という動力学的帰結の、構造論的根拠である。$\mathcal{R}$ の強化は $\mathcal{T}$ を生まない ── なぜなら $\mathcal{T}$ は $\mathcal{R}$ の \textbf{強化の方向} には存在しないからである。

\subsection{Failure modes of the individual node}


\subsubsection{The structural relation to \S{}7.4}


\S{}7.4 で示した個体内現象 ── explanatory overreach, compulsive coherence-seeking, illusions of global optimization, premature or forced conclusions ── は、本セクションの構造論的事実から \textbf{同じ場所} に到達する。

両セクションの関係:

\begin{itemize}
\item \S{}7.4 は \textbf{動力学的視点} から症状を記述した(正しさの累積による分岐爆発の帰結として)
\item \S{}8.4 は \textbf{構造論的視点} から同じ症状を記述する(停止演算子の不在の帰結として)
\end{itemize}

これらは同じ現象の二つの記述である。\S{}7 の動力学と \S{}8 の構造は、互いに独立した二つの問題ではなく、同じ構造的事実の異なる側面を捉えている。

\subsubsection{The structural reading of the symptoms}


\S{}7.4 で挙げた症状を、\S{}8 の構造論的視点から読み直す:

\textbf{Over-explanation(過剰説明):}
停止演算子が存在しないため、説明連鎖は内発的に終端を持てない。個体は連鎖を延長し続けるしかなく、その延長は時に説明されるべき対象の範囲を超えて続く。これは説明能力の過剰ではなく、停止能力の不在である。

\textbf{Compulsive coherence-seeking(強迫的整合性追求):}
停止演算子が存在しないため、矛盾検出 → 矛盾解消 → 新たな矛盾検出という連鎖が終端を持てない。個体は整合性を求め続けるしかなく、整合性そのものが追求の終端を提供する構造的根拠を持たない。

\textbf{Illusions of total understanding(全的理解の錯覚):}
停止演算子が存在しないため、推論連鎖の途中で「ここまでで十分」という判定を内発的に行えない。代わりに個体は、現在の到達点が \textbf{全体に到達した点} であるかのように錯覚することで、機能的停止を擬似的に達成する。これは判断ではなく構造的代替物である。

\textbf{Forced or premature conclusions(強引な結論):}
停止演算子が存在しないにもかかわらず、何らかの結論が要請される状況では、個体は構造的根拠を欠いた閉包を強引に遂行する。これは個体の判断力不足ではなく、停止が要請されるが停止演算子が存在しないという構造的矛盾の現れである。

\subsubsection{Not pathologies in the medical sense}


\S{}7.4.2 で強調した通り、これらの症状は医学的・心理学的意味での病理ではない。本セクションの構造論的視点からは、これらは:

\begin{itemize}
\item 個体系という構造的単位の \textbf{構造的限界} の自然な現れ
\item 停止演算子の不在から動的に生成される \textbf{典型的応答パターン}
\item 「異常」ではなく「個体内に閉じた知性的動態の標準的振る舞い」
\end{itemize}

である。これらが「症状」として現れるのは、外部からそれを観察する場合や、個体自身が外的基準と照合する場合のみである。個体内部からは、これらは単に「推論を続けている」状態として現れる。

\subsection{Anticipated objection: stopping versus interruption}


\subsubsection{The objection}


\S{}8.3 の主張 ── 個体内には停止演算子が構造的に存在しえない ── に対して、直感的な反論が想定される:

\begin{quote}
実際の人間は、因果推論を無限に続けない。疲労によって思考を止め、報酬勾配の減衰によって興味を失い、注意の漂流によって他の対象に移り、生理的要請(睡眠、空腹、痛み)によって推論連鎖を打ち切る。これらはすべて個体内で起こっている。ゆえに個体内に停止機構は存在する。
\end{quote}

この反論は直感的には強力であり、適切に応答しなければ \S{}8 の中核論証が崩される。

\subsubsection{The structural distinction: stopping versus interruption}


応答の核心は、\textbf{停止(stopping)} と \textbf{中断(interruption)} の構造的区別にある。

\textbf{停止(stopping)} とは:
\begin{itemize}
\item 推論連鎖がその内部から「完了した」「これ以上は不要である」と判定されること
\item 連鎖の内的構造によって終端が確定されること
\item 同じ連鎖の再開を構造的に必要としない状態
\end{itemize}

\textbf{中断(interruption)} とは:
\begin{itemize}
\item 推論連鎖が外的要因によって途切れること
\item 連鎖の内的構造とは独立に、外的条件によって動作が止まること
\item 外的条件が変化すれば連鎖が再開されうる状態
\end{itemize}

両者は表面上の現れ ── 「推論が止まる」 ── は類似するが、構造的にはまったく異なる事象である。

\S{}8.3 で論じた「停止演算子の構造的不在」は、停止についての主張であって、中断についての主張ではない。中断は構造的に常に可能である ── ただし中断は連鎖の \textbf{完了} を意味せず、外的条件が変われば連鎖は再開される。

\subsubsection{Energy, reward gradients, and biological constraints}


反論で挙げられた要因 ── 疲労、報酬勾配の減衰、注意の漂流、生理的要請 ── を構造論的に位置づける。

これらはいずれも \textbf{生命層(Bio-IFGT)の substrate-level な制約} である。具体的には:

\begin{itemize}
\item \textbf{エネルギー枯渇} は代謝制約であり、生命層の準閉包 regime における動作可能エネルギー範囲の限界を示す
\item \textbf{報酬勾配の減衰} は神経学的・生化学的応答の動的特性であり、生命層のフィードバック構造に属する
\item \textbf{注意の漂流} は神経活動の競合的動的構造であり、生命層の状態切り替え機構に属する
\item \textbf{生理的要請} は代謝・恒常性維持の周期的要求であり、生命層の維持機構に属する
\end{itemize}

これらはいずれも因果推論の構造そのものには属さない。それらは因果推論を \textbf{担う substrate} に属する。生命層がこれらの制約を持つことは、知性層に停止演算子があることを意味しない ── 両者は \S{}9.4 で示される異なる構造層に属する現象である。

具体的に確認すれば、これらの要因による「停止」は実際には \textbf{中断} である:

\begin{itemize}
\item 疲労によって止まった推論は、休息後に同じ問いから再開される(Gettier 問題、Fermat の最終定理、未解決の人生上の問い ── これらは数十年にわたって個体的に中断と再開を繰り返してきた)
\item 報酬勾配の減衰で離れた問いは、状況が変われば再び興味の対象となる
\item 注意の漂流で逸れた推論は、合図によって呼び戻されうる
\item 生理的中断は意図的訓練(瞑想、専門化された訓練)によって部分的に克服される
\end{itemize}

これらすべてにおいて、中断は連鎖の \textbf{完了} ではない。連鎖は依然として停止演算子を欠いており、外的要因による中断と再開を繰り返しているにすぎない。

\subsubsection{Why this distinction is decisive: the paperclip problem}


停止と中断の混同は単なる用語論の問題ではない。それは AI 安全性議論において六十年近く解決を見ない核心的問題 ── ペーパークリップ問題 ── の構造的根源である。

Bostrom 由来のペーパークリップ最大化思考実験は、目的関数を持つ強力な AI が、その目的の追求を停止する内発的根拠を持たないがゆえに、最終的に宇宙的規模の最適化を遂行しうることを示す。これに対する伝統的応答は、本論文の枠組みからは構造的に誤っている:

\textbf{「正しい目的を設計すればよい」という応答} は、\S{}7 で示した正しさの構造を見落としている。目的の正しさは停止を生まない。正しい目的は、むしろ再利用可能性と推論連結性によって、分岐爆発を加速する。「アラインメント問題」を停止条件の精緻化として解こうとする限り、これは \S{}6 で論じた Gettier 問題の修正試みと同じ構造的失敗を繰り返す。

\textbf{「十分に賢い AI なら自分で停止する」という応答} は、\S{}7.5 で示した「より多くの知性は問題を加速する」という構造的事実を見落としている。知性的動態の強化は停止演算子を生まない。

本論文の枠組みからのペーパークリップ問題の構造論的同定は以下である:

\begin{quote}
ペーパークリップ問題は AI 固有の問題ではない。
それは個体内停止演算子の構造的不在 ── すなわち知性層における local horizon-form の一般的問題 ── が、
生命層の中断機構を持たない人工系において直接顕現する場所である。
\end{quote}

人間が「ペーパークリップ問題を起こさないように見える」のは、人間に停止演算子があるからではない。人間が生命層の制約(エネルギー、報酬勾配、生理的要請)による \textbf{中断} を絶え間なく受け続けているからである。これらの中断が、表面的には「自然な停止」と見えるパターンを生む。

しかし \S{}5.2 で扱った Transformer のような人工系は、これら生命層の中断機構を持たない。計算資源が許す限り、推論は分岐爆発のレジームに留まり続ける。ここでは知性層の構造的不在が、生命層によって緩衝されることなく直接現れる。

ゆえに、ペーパークリップ問題への構造的応答は \textbf{アラインメント} でも \textbf{能力強化} でもなく、\textbf{越境の構造化} ── すなわち \S{}9 で論じる local horizon-form を、人工系においても明示的に設計可能な外化構造として組み込むこと ── である。これは本論文の射程外の課題であり、後続論文(Part II)で扱われる。

\subsubsection{What the objection does not refute}


以上から、反論で挙げられた要因は \S{}8 の中核主張を反駁しない。むしろ、これらの要因は \S{}8 の主張を \textbf{間接的に支持する}:

\begin{itemize}
\item 因果推論が個体内で「自然に止まる」ように見える現象は、すべて生命層の中断機構によって説明される
\item 生命層の中断機構を欠いた人工系では、知性層の構造的不在が \textbf{より純粋な形で現れる}(ペーパークリップ問題)
\item 中断と停止を混同しない限り、\S{}8 の構造論的主張は維持される
\end{itemize}

したがって反論への応答は、\S{}8 の論証を修正することではなく、\textbf{何を「停止」と呼ぶかの構造的明確化} である。この明確化のもとで、\S{}8 の主張は次のように再述される:

\begin{quote}
個体内には、推論連鎖を \textbf{完了} させる演算子が構造的に存在しえない。
個体内で観察される推論の終了は、すべて中断であって停止ではない。
真の意味での停止は、\S{}9 で論じる越境境界の構造を経由してのみ可能である。
\end{quote}

\subsection{Implications for the broader argument}


\subsubsection{The transition to \S{}9}


\S{}8 で示した構造的事実 ── 個体内には停止演算子が構造的に存在しえない ── は、\S{}9 で展開される議論への直接的入口となる。

停止が個体内では生成されない以上、停止は \textbf{個体の外側} にある記録・応答・合意・遅延の構造を要請する。これは選択ではなく構造的制約である。\S{}9 ではこの越境点を概念化し、命名する:

\begin{quote}
個体内では到達できず、越境することによってのみ推論の停止と持続が可能になる構造的境界を、
\textbf{the differential horizon of intelligence} と呼ぶ。
\end{quote}

\S{}9 で示される越境境界は、\S{}8 で論証した構造的不在を \textbf{裏返した側面} である。何が個体内にないかを \S{}8 が示し、何が個体外に要請されるかを \S{}9 が示す。両者は同じ構造的事実の異なる視点である。

\subsubsection{Connection to the layered correspondence (\S{}9.4)}


\S{}9.4 で示される物理層・生命層・意識層・知性層の四層対応のうち、知性層における local horizon-form は本セクションの構造論的事実から導かれる。物理層・生命層・意識層との対応は、各層において同様の \textbf{個体的閉包の不可能性} が異なる仕方で現れていることを示す。

特に \S{}8.5 で論じたエネルギー・報酬・代謝制約の生命層への位置づけは、\S{}9.4 における四層対応において、生命層の constituent な特徴として再び現れる。また、個体内における temporal log referencing(意識編の主題)は、\S{}9.4 で意識層として、生命層と知性層の遷移層として位置づけられる。

すなわち本セクションは、知性層における local horizon-form の \textbf{構造的根拠} を提供する位置にある。

\subsection{Summary}


本セクションの結論を以下に述べる:

\begin{quote}
個体内には、知性的動態を停止させる演算子が \textbf{構造的に存在しえない}。
これは個体の欠陥や能力不足から生じるのではなく、
因果推論の再帰構造そのものから導出される構造的事実である。
\end{quote}

\begin{quote}
停止のための四条件 ── 外部応答、共有合意、log としての外化、時間的遅延 ── は
いずれも個体的推論の外側に位置する。
個体的推論からは、これら条件のいずれも内発的に生成されない。
\end{quote}

\begin{quote}
個体的知性は、誤っているから失敗するのでも、能力不足から失敗するのでもない。
停止できないから失敗する。
これは欠陥ではなく構造的欠在(structural absence)である。
\end{quote}

\begin{quote}
個体内で観察される推論の終了は、すべて中断であって停止ではない。
中断は生命層の制約(エネルギー、報酬勾配、生理的要請)によって与えられ、
知性層に停止演算子があることを意味しない。
この区別の欠如が、ペーパークリップ問題などの AI 安全性議論を構造的に停滞させてきた。
\end{quote}

この構造的欠在は、\S{}9 で展開される越境境界の必然性の根拠となる:

\begin{itemize}
\item \S{}9 では、この構造的欠在を \textbf{裏返した側面} として越境境界が命名される
\item \S{}9.4 では、本セクションの構造的事実が物理層・生命層・意識層と対応する四層構造として位置づけられる
\item \S{}10 では、本セクションが示した構造的事実が論文全体の制約の一つとして集約される
\end{itemize}
